\subsection{GERMS AT A POINT}

One of the commonest situations to which the definitions and results of no. 9 apply is that in which \(\mathfrak{F}\) is the neighbourhood filter of a point \(a\) of a topological space \(\mathbf{X}\) ; instead of ``germs with respect to \(\mathfrak{F}\) '' we then speak of ``germs at the point \(a\) ''. Notice that there is only one germ of neighbourhoods of \(a\) , namely the germ of the whole space \(\mathbf{X}\) . The germs of closed sets are identical with the germs of sets which are locally closed at the point \(a\) , for if \(\mathbf{L}\) is locally closed at \(a\) , then the germs of \(\mathbf{L}\) and \(\overline{\mathbf{L}}\) at \(a\) are equal ( \(\S 3\) , no. 1, Proposition 1). It follows that if \(\xi, \eta\) are two germs of locally closed sets at \(a\) , then so are \(\xi \cup \eta\) and \(\xi \cap \eta\) .

Since \(a\) is in each \(\mathbf{V} \in \mathfrak{F}\) , \(f(a)\) is defined for each mapping \(f\) whose domain belongs to \(\mathfrak{F}\) ; furthermore, if \(f\) and \(g\) have the same germ at \(a\) we must have \(f(a) = g(a)\) , so that \(f(a)\) depends only on the germ \(\tilde{f}\) of \(f\) at \(a\) , and is called the value of \(\tilde{f}\) at \(a\) and is denoted by \(\tilde{f}(a)\) . It should be emphasized that the relation \(\tilde{f}(a) = \tilde{g}(a)\) does not in general imply that \(\tilde{f} = \tilde{g}\) .

Let \(X', X''\) be two topological spaces; \(b\) a point of \(X''\) ; \(g, g'\) two mappings of \(X'\) into \(X''\) having the same germ at \(b\) . If \(f, f'\) are two mappings of \(X\) into \(X'\) which are continuous at \(a\) and have the same germ at \(a\) and are such that \(f(a) = b\) , then \(g \circ f\) and \(g' \circ f'\) have the same germ at the point \(a\) ; for if \(V'\) is a neighbourhood of \(b\) in \(X'\) such that \(g(x') = g'(x')\) for all \(x' \in V'\) , then there is a neighbourhood \(V\) of \(a\) such that \(f(V) \subset V'\) , \(f'(V) \subset V'\) and \(f(x) = f'(x)\) for all \(x \in V\) , and the assertion follows. The germ of \(g \circ f\) at \(a\) is then called the composition of the germs \(\tilde{g}\) and \(\tilde{f}\) of \(g\) and \(f\) respectively and is denoted by \(\tilde{g} \circ \tilde{f}\) .

